%! Author = Alek
%! Date = 3/15/2026

\documentclass{article}
\usepackage{graphicx} % Required for inserting images
\usepackage{amsmath}
\usepackage[margin=2cm]{geometry}
\usepackage{gensymb}
\usepackage{float}




\title{Multi-Variable Control Systems Project Update 1}
\author{Alek Krupka}
\date{February 2026}

\begin{document}
    \maketitle

    \section{Introduction}\label{sec:introduction}

    For this project update we derive the 4-block representation of the control problem along with showing the MIMO
    system representation from Assignment 3.

    \section{4-Block Representation}\label{sec:four-block-representation}

    To find the 4 block-representation we need to select our objective to minimize.
    Our 4-block representation will be done on our SISO system from the first project update.

    We want to select P to minimize the mixed sensitivity function shown below.

    \begin{align}
        \min_K \left\| \begin{array}{c} W_y T_o  W_n \\ W_u K S_o W_n \\ W_e S_o W_r \end{array} \right\|_{H2}
    \end{align}

    To do this, we select our four-block-representation to be the following.

    \begin{align}
        P = \begin{bmatrix}
                \begin{bmatrix} 0 & 0 & 0 \\ 0 & 0 & 0 \\ 0 & 0 & W_e W_r \end{bmatrix} &
                \begin{bmatrix} W_y G \\ W_u \\ -W_e \end{bmatrix} \\
                \begin{bmatrix} W_n & W_n & W_r \end{bmatrix} & -G
        \end{bmatrix}
    \end{align}

    To find $G$ we just use the formula $G = C(sI-A^{-1})B$.
    The weightings other transfer functions are tuned and can be seen in the MATLAB live
    script.

    \section{H2 Norm Gain Calculation}\label{sec:H2_Norm}

    To calculate the H2 norm of the calculation we call the MATLAB function "h2syn" which
    designs the controller $K$ to control the plant.

    See the MATLAB live code for the controller $K$ of for the closed loop systems.

    \section{H-$\infty$ Norm Gain Calculation}\label{sec:Hinf_Norm}

    To calculate the H2 norm of the calculation we call the MATLAB function "h2infsyn" which
    designs the controller $K$ to control the plant.

    See the MATLAB live code for the controller $K$ of for the closed loop systems.

    \section{MIMO System Design}\label{sec:MIMO_Design}

    For this assignment, we will be doing analysis on the MIMO drone control system.
    To do this, we will of course need to linearize the system around a specific operating point.
    We select this point to be equilibrium (where the drone is not rotating or accelerating in any direction).
    In this case, we want to find the configuration for which $\alpha_y$ and $a_z = 0$ and $a_x = 0$.
    \\
    To do this, we can use our work done previously to set this operating point to the following,

    \begin{align}
        \omega_1 = \omega_2 = 495 rad/s
    \end{align}

    Under the same model assumptions as before with

    \[ k_f = 10^{-5} \pm 10^{-6}, m = 0.5 \pm 0.1kg, I_{rod} \approx \frac{2}{3} \times 10^{-4}, l = 0.1m \pm 5mm\]

    Using these variables we arrive at the equations,

    \begin{align}
        a_z &= 2\times 10^{-5} \cos(\theta) (\omega_0^2 + \omega_1^2)-9.81\\
        a_x &= 2\times 10^{-5} \sin(\theta) (\omega_0^2 + \omega_1^2)\\
        \dot{v}_z &= a_z\\
        \dot{v}_x &= a_x\\
        \dot{z} &= v_z\\
        \dot{x} &= v_x\\
        \dot{\omega_y} &= 1.5 \times 10^{-2}(\omega_1^2 - \omega_2^2)\\
        \dot{\theta} &= \omega_y
    \end{align}

    Linearizing our squared terms, we get the that.

    \begin{align}
        \omega_i^2 &= \omega^{\star} + 2\Delta \omega_i\\
        &= 495^2 + 2 \Delta \omega_i
    \end{align}

    For $a_z$ the $\omega^{\star}$ cancels out with gravity since we linearize around $\theta = 0$.
    After cancelling, we see that.

    \begin{align}
        a_z = 4 \times 10^{-5} (\Delta \omega_1 + \Delta \omega_2)
    \end{align}

    However, for finding $a_x$ and linearizing $\sin(\theta)$ around 0, we find that.

    \begin{align}
        a_x &= 2 \times 10^{-5} \theta((\omega_i^{\star} + \Delta \omega_i)^2 + (\omega_i^{\star} + \Delta \omega_i)^2)\\
    \end{align}

    Keeping only the linear terms (ie.
    terms that don't have the input squared or input multiplied by $\theta$) we see that.

    \begin{align}
        a_x &= 4 \times 10^{-5} * 495^2 \theta
    \end{align}

    Finally, linearizing $\omega_y$ yields the resulting equation.

    \begin{align}
        \alpha_y &= 0.03(\Delta \omega_1 - \Delta \omega_2)
    \end{align}

    Using the newly linearized equations, we can represent the state space
    as shown below.

    \begin{align}
        \dot{x} &= \begin{bmatrix}
                       0 & 1 & 0 & 0 & 0 & 0 \\
                       0 & 0 & 0 & 0 & 9.801 & 0 \\
                       0 & 0 & 0 & 1 & 0 & 0 \\
                       0 & 0 & 0 & 0 & 0 & 0 \\
                       0 & 0 & 0 & 0 & 0 & 1 \\
                       0 & 0 & 0 & 0 & 0 & 0 \\
        \end{bmatrix}
        x + \begin{bmatrix} 0 & 0 \\ 0 & 0 \\ 0 & 0 \\ 4 \times 10^{-5} &
        4 \times 10^{-5} \\ 0 & 0 \\ 0.03 & -0.03 \end{bmatrix}u\\
        y &= \begin{bmatrix} 1 & 0 & 0 & 0 & 0 & 0 \\
            0 & 0 & 1 & 0 & 0 & 0 \\ 0 & 0 & 0 & 0 & 1 & 0 \end{bmatrix}x
    \end{align}

    Using this state space representation, we see that.

    \begin{align}
        A &= \begin{bmatrix}
                 0 & 1 & 0 & 0 & 0 & 0 \\
                 0 & 0 & 0 & 0 & 9.801 & 0 \\
                 0 & 0 & 0 & 1 & 0 & 0 \\
                 0 & 0 & 0 & 0 & 0 & 0 \\
                 0 & 0 & 0 & 0 & 0 & 1 \\
                 0 & 0 & 0 & 0 & 0 & 0 \\
        \end{bmatrix}\\
        B &= \begin{bmatrix} 0 & 0 \\ 0 & 0 \\ 0 & 0 \\ 4 \times 10^{-5} &
        4 \times 10^{-5} \\ 0 & 0 \\ 0.03 & -0.03 \end{bmatrix}\\
        C &= \begin{bmatrix} 1 & 0 & 0 & 0 & 0 & 0 \\
            0 & 0 & 1 & 0 & 0 & 0 \\ 0 & 0 & 0 & 0 & 1 & 0 \end{bmatrix}
    \end{align}

    This gives our output in terms of $y = \begin{bmatrix} x, z, \theta \end{bmatrix}^T$.\\
    \\
    Next, to have a stable system, we apply full state feedback to the system.\\
    While the final system will need an observer as we cannot add sensors for every state, that observer will be designed at a later date.
    In this case, though, we use MATLAB to define the poles for our system at $-1 \pm 0.1j$ and 4 poles at $-2 \pm 0.1j$.\\
    \\
    This will allow us later to find the controllability and observability Gramians of our MIMO system.

    \section{Transmission Poles and Zeros}\label{sec:tzeros_and_poles}
    Using these matrices, we can find the transmission poles and zeros of this state space.
    Define the open loop transfer function $G(s)$ as follows.

    \begin{align}
        G(s) &= C(sI-A)^{-1}B
    \end{align}

    Now, we need to find the smith form of the state space representation to find the transmission poles and zeros.
    To do this, we can use MATLAB though.\\
    \\
    However, MATLAB will only turn square transfer functions into Smith McMillan form.
    As a result, our code just defines the system and finds the transmission zeros and poles via the "pole" and "tzero" function.\\
    \\
    Looking at the MATLAB script in the appendix, we see that our transmission poles and zeros are the following.\\
    \\
    \begin{itemize}
        \item Poles: $-2 \pm 0.1j$, $-1 \pm 0.1j$ each with geometric multiplicity of 2.
        This totals 6 poles and are exactly as we designed.
        \item Transmission Zeros: The system has no transmission zeros as calculated by MATLAB.
    \end{itemize}

    \section{Controllability and Observability Grammians}\label{sec:controllability_and_observability}

    According to MATLAB (as shown by the attached script), the controllability and observability Gramians are as follows.

    \subsection{Controllability Gramian}\label{subsec:Controllability Gramian}

    As found by MATLAB the controllability Gramian for our system is the following,

    \begin{align}
        W_c &= 10^{-3}\begin{bmatrix} 2.2 & 0 & -6 & 1.9 & -0.1 & 0 \\ 0 & 0.8 & -1.9 & -2.2 & 0 & -0.2 \\ -6 & -1.9 & 20.9 & 0 & 0.2 & 0.4 \\ 1.9 & -2.2 & 0 & 15.1 & -0.4 & 0.4 \\ -0.1 & 0 & 0.2 & -0.4 & 0 & 0 \\ 0 & -0.2 & 0.4 & 0.4 & 0 & 0.2 \end{bmatrix}
    \end{align}

    The eigenvalues found by MATLAB are the following.

    The eigenvalues given by the controllability matrix are the following vector.

    \begin{align}
        10^{-3} \begin{bmatrix} 0 & 0 & 0.1 & 0.5 & 15.7 & 22.8 \end{bmatrix}
    \end{align}

    As we can see, the smallest eigenvalue generated from our controllability gramian is 0 indicating that we have states that cannot be controlled easily.
    The largest eigenvalue is 0.22 indicating the system is not very controllable overall.
    The first two states being mostly uncontrollable make sense since in order to control the x-position
    we need to control the angle of the drone, which is difficult to do by itself.
    Even then, at a small angle, the x-position will hardly change.

    \subsection{Observability Gramian}\label{subsec:observability_gramian}

    A found by MATLAB the observability Gramian of our closed loop system is the following.

    \begin{align}
        W_o &= \begin{bmatrix} 4.6324 & 6.6149 & -0.6419 & -0.2019 & 35.4426 & 6.6023 \\ 6.6149 & 11.0248 & -1.2017 & -0.3781 & 61.6527 & 11.6924 \\
            -0.6419 & -1.2017 & 0.6238 & 0.1247 & -7.0574 & -1.3722 \\ -0.2019 & -0.3781 & 0.1247 & 0.0312 & -2.2217 & -0.4321 \\
            35.4426 & 61.6527 & -7.0574 & -2.2217 & 350.6472 & 66.9726 \\
            6.6023 & 11.6924 & -1.3722 & 0.4321 & 66.9726 & 12.8418
        \end{bmatrix}
    \end{align}

    The eigenvalues for the controllability matrix are shown below.

    \begin{align}
        10^{-3} \begin{bmatrix} 3.4 \\ 9.2 \\ 50 \\ 484.9 \\ 1217.9 \\ 378000 \end{bmatrix}
    \end{align}

    As we can see, the smallest eigenvalues of our system is roughly $3 \times 10^{-3}$ while our largest
    is over 300.
    As a result, the difference between our most observable and least observable states is massive and will prove a challenge for observing the state of the drone.\\
    \\
    However, this result does make sense as by our original equations, any small difference is motor speed will result in a large change of angle
    meaning the state is very observable.

    \section{Balanced Realization}\label{sec:balanced_realization}

    Next we will find the balanced realization of the system.
    This is the modified state-space where both the controllability Gramians are diagonal matrices that are equivalent.
    This space can again be found using MATLAB and is shown in the attachment.
    The balanced realization from $u$ to $y$ is also shown below.

    Let $\hat{A}, \hat{B}, \hat{C}, \hat{D}$, be the balanced state space realization.

    \begin{align}
        \hat{A} &= \begin{bmatrix} -0.2232 & 0.6697 & 0.01046 & 0.1821 & 7.929 \times 10^{-4} & -8.132 \times 10^{-4} \\
            -0.6755 & -1.045 & -0.549 & -0.7724 & -0.004782 & 0.003034 \\
            0.2361 & 1.131 & -0.9512 & 0.7909 & -0.01922 & 0.002614 \\
            0.2065 & 0.9147 & -1.833 & -3.78 & 0.00832 & -0.001142 \\
            8.534 \times 10^{-7} & 3.026 \times 10^{-6} & 8.182 \times 10^{-7} & 1.89 \times 10^{-7} & -0.7792 & 0.8828 \\
            1.097 \times 10^{-6} & 3.857 \times 10^{-6} & -3.897 \times 10^{-6} & -3.816 \times 10^{-6} & -1.7 & -3.221 \end{bmatrix}\\
        \hat{B} &= \begin{bmatrix} -0.209 & 0.2091 \\ -0.233 & 2.33 \\ 0.1109 & -0.1108 \\ 0.09781 & -0.09774 \\ -0.001158 & -0.001159 \\ -0.001116 & -0.001117 \end{bmatrix}\\
        \hat{C} &= \begin{bmatrix} -0.07062 & 0.13 & 0.1553 & -0.01723 & 0.001607 & -0.0002296 \\
            0.2871 & -0.3025 & 0.0206 & -0.1303 & -0.0001449 & 0.0004956 \\
            -0.0004693 & -0.01493 & 0.005565 & -0.04289 & 0.0002854 & -0.001481 \end{bmatrix}\\
        \hat{D} &= 0
    \end{align}

    The henkel singular values for the balanced realization were found using MATLAB and are listed \textbf{below}

    \begin{align}
        \sigma = \begin{bmatrix} 0.1957 \\ 0.0520 \\ 0.0129 \\ 0.0025 \\ 0 \\ 0 \end{bmatrix}
    \end{align}

    Again, this makes sense as two of the states are not very controllable.
    The most controllable and observable state of the transformed state space is the first state
    as it has the largest eigenvalues.
    Meanwhile the last state is the least controllable as it has eigenvalues of 0 in the transformed state space.

    \section{Appendix}\label{sec:appendix}

    See the attached MATLAB live script for all calculations in MATLAB.


\end{document}

